\documentclass[a4paper]{amsart}
\usepackage[a4paper,margin=1in]{geometry}
\usepackage{amsmath,amssymb,amsthm}
\usepackage{booktabs,array}

%--- theorem environments ---------------------------------------------------
\newtheorem{theorem}{Theorem}
\newtheorem{corollary}[theorem]{Corollary}
\newtheorem{remark}[theorem]{Remark}
\theoremstyle{definition}
\newtheorem{assumption}{Assumption}

\begin{document}

\title{Travelling-Wave Theory for CO$_2$ Adsorption Breakthrough\\
with Dual-Kinetic Amine Kinetics and Toth Equilibrium}

\date{}
\maketitle

%--- Abstract ---------------------------------------------------------------
\begin{abstract}
Direct air capture of CO$_2$ at ambient concentrations near 420~ppm requires
packed-bed amine sorbents whose breakthrough timing and curve shape depend
jointly on equilibrium heterogeneity and kinetic multiplicity.
For polyethylenimine-impregnated silica (PEI@SiO$_2$), two structurally
distinct amine populations---fast surface amines and slow bulk-polymer
amines---produce a sharp sigmoidal front followed by a prolonged tail that
single-timescale kinetic models cannot reproduce.
We derive, from the one-dimensional axially-dispersed species balance with a
temperature-dependent Toth closure, a travelling-wave reduction valid under a
three-level hierarchy of small parameters that holds rigorously at
DAC-relevant conditions.
Theorem~\ref{thm:wave} proves that the dual-kinetic LDF system admits a
travelling wave whose speed is set by the fast-site equilibrium capacity
alone, and that the slow-site contribution is a uniformly bounded
$O(\kappa)$ correction decaying asymptotically in the wave frame; the
resulting hybrid outlet expression reduces exactly to the Myers--Font (2020)
single-LDF result when the sorbent is kinetically homogeneous.
A scalar asymmetry index $\mathcal{A}=(\tau_{95}-\tau_{50})/(\tau_{50}-\tau_5)$
converts any measured breakthrough curve into a kinetic-model selection
criterion, with $\mathcal{A}\lesssim 2$ selecting the single-LDF closure
and $\mathcal{A}\gtrsim 3$ triggering the dual-kinetic extension.
Three escalating validation gates---Klinkenberg linear benchmark, Rankine--Hugoniot
front velocity to within $\pm15\%$, and Stampi-Bombelli et al.\ (2024)
breakthrough time to within $\pm20\%$---constitute a prior-independent
commitment protocol before any parametric sweep.
\end{abstract}

%============================================================================
\section{Introduction}
%============================================================================

Atmospheric CO$_2$ exceeded 422~ppm in 2024.  Direct air capture (DAC) has
entered techno-economic roadmaps as a necessary complement to source-side
reductions in hard-to-decarbonise sectors \cite{IEA2022,IPCC2023}.  Among
scalable DAC hardware, packed-bed columns loaded with amine-functionalised
solid sorbents offer high volumetric capacity, operational flexibility, and
compatibility with temperature-swing regeneration; their process economics have
recently been quantified for PEI@SiO$_2$ granules in reactive-transport cycle
models \cite{Pedrozo2025}.

The central engineering question is: given a sorbent with a characterised
equilibrium isotherm and a known mass-transfer coefficient, when does CO$_2$
break through the column, and what shape does the outlet concentration profile
take as a function of the design variables---superficial velocity $u$, inlet
concentration $c_0$, bed length $L$, and temperature $T$?  The breakthrough
time $\tau_\text{BT}$ and the curve shape together determine bed utilisation,
dynamic capacity, and cycle energy penalty.

Existing parametric studies operate chiefly at post-combustion concentrations
($y_{\mathrm{CO}_2}\sim5$--15\%) where gas-phase accumulation is non-negligible.
At DAC concentrations, the inlet CO$_2$ mole fraction is
$y_0 = c_0 R_g T/p_a \approx 4\times10^{-4}$, and the solid stores
$\Lambda = (1-\varepsilon)\rho_p q_0^*/(\varepsilon c_0)\sim10^{4}$--$10^{5}$
times more CO$_2$ than the gas phase.  This hierarchy collapses the governing
partial differential system to a one-parameter family of ODEs---the
constant-pattern or travelling-wave reduction---whose mathematical properties
in the DAC limit have not been systematically exploited.

A second gap is kinetic.  PEI@SiO$_2$ exhibits two structurally distinct amine
populations: surface-accessible primary amines that react via carbamate
formation on timescales of minutes, and bulk polymer-layer amines that require
diffusion through the viscous polyethylenimine matrix before reacting, with
timescales one to two orders of magnitude longer \cite{Bollini2012,%
Kalyanaraman2015,Ohs2018}.  The resulting breakthrough curve carries a sharp
initial rise followed by a slowly decaying tail---a morphology that a
single-timescale linear driving force (LDF) model cannot reproduce and that
empirical models (Bohart--Adams, Thomas, Yoon--Nelson) cannot physically
explain, because all three reduce to the same logistic function with drifting
parameters \cite{Hu2024}.

This paper addresses both gaps.  We identify three small parameters
$(\delta_1, 1/\mathrm{Pe}, \kappa)$ natural to DAC operating conditions, carry out
a systematic asymptotic reduction of the coupled PDE--LDF system to a
travelling-wave ODE, and prove that the dual-kinetic extension of this wave
is well-posed.  The resulting hybrid analytical expression is exact in the
limits of homogeneous kinetics, DAC incompressibility, and favourable Toth
isotherm, and reduces to the Myers--Font (2020) result \cite{MyersFont2020}
when $\eta\to1$.

We also introduce a scalar asymmetry index that turns a shape measurement
into a kinetic-model selection criterion, and we define three escalating
validation gates---rather than a single comparison---as a commitment device
against overfitting.  Several notation conflicts present in the literature are
resolved, including the form of the LDF driving force
(Section~\ref{sec:resolution}) and the definition of the gas-to-solid
capacity ratio $\delta_1$ (Table~\ref{tab:dimless}).

%============================================================================
\section{Governing equations}
%============================================================================

\subsection{Physical system and assumptions}

A cylindrical packed bed of length $L$ and inner diameter $d_c$ is filled
with PEI@SiO$_2$ granules of mean diameter $d_p$ and pellet density $\rho_p$.
A binary CO$_2$/N$_2$ stream enters at $z=0$ with superficial velocity $u$,
CO$_2$ concentration $c_0$, and temperature $T_0$; the bed is initially
clean.

\begin{assumption}[Plug flow]\label{ass:plug}
No radial gradients; the particle P\'eclet number $Pe_p = ud_p/D_m > 10$
justifies neglecting radial dispersion.
\end{assumption}

\begin{assumption}[Axial dispersion]\label{ass:ax}
Axial backmixing is retained via a lumped coefficient $D_L$ correlated to
velocity and particle size \cite{Stampi2024}.
\end{assumption}

\begin{assumption}[Ideal gas]\label{ass:ig}
$p = cR_gT$, where $c$ is total molar concentration.
\end{assumption}

\begin{assumption}[Solid-loading LDF]\label{ass:ldf}
The solid uptake rate is proportional to the solid-loading driving force:
$\partial\bar{q}/\partial t = k(q^* - \bar{q})$, where $q^*(c,T)$ is the
Toth equilibrium loading at local conditions.  This form---not the
gas-concentration form $k'(c-c^*)$---is used throughout; the two forms use
different rate coefficients and must not be conflated.
\end{assumption}

\begin{assumption}[Dual kinetics]\label{ass:dk}
The total loading $\bar{q}=q_1+q_2$ splits into a fast-site component $q_1$
(surface amines, equilibrium fraction $\eta$, rate constant $k_1$) and a
slow-site component $q_2$ (bulk-polymer amines, equilibrium fraction $1-\eta$,
rate constant $k_2 = \kappa k_1$) with $0 < \kappa \ll 1$.
\end{assumption}

\begin{assumption}[Isothermal screening]\label{ass:iso}
For the analytical reduction, $T = T_0$ is fixed.  Non-isothermal effects
enter at relative order $\Phi\cdot\delta_1$ (Table~\ref{tab:dimless}) and
are retained in the full numerical model via the energy balance
\eqref{eq:energy}.
\end{assumption}

\subsection{Species mass balance}

The CO$_2$ gas-phase balance is
\begin{equation}
  \varepsilon\frac{\partial c}{\partial t}
  + u\frac{\partial c}{\partial z}
  = \varepsilon D_L\frac{\partial^2 c}{\partial z^2}
    - (1-\varepsilon)\rho_p\frac{\partial\bar{q}}{\partial t},
  \label{eq:gas}
\end{equation}
where $\varepsilon$ is the \emph{bed void fraction}, distinct from the particle
porosity $\varepsilon_p$ (for $\gamma$-alumina pellets, $\varepsilon_p\approx0.71$
\cite{Stampi2024}; the two quantities appear in separate correlations and must
not be conflated).  The dual-kinetic LDF closures are
\begin{align}
  \frac{\partial q_1}{\partial t}
    &= k_1\bigl(\eta\,q^*(c,T_0) - q_1\bigr),
  \label{eq:q1}\\
  \frac{\partial q_2}{\partial t}
    &= \kappa k_1\bigl((1-\eta)\,q^*(c,T_0) - q_2\bigr),
  \label{eq:q2}
\end{align}
with $\bar{q}=q_1+q_2$.

\subsection{Toth equilibrium isotherm}

The temperature-dependent Toth isotherm \cite{Grossmann2023,Stampi2024} is
\begin{equation}
  q^*(p,T) = \frac{n_s(T)\,b(T)\,p}
               {\bigl[1+(b(T)\,p)^{t(T)}\bigr]^{1/t(T)}},
  \label{eq:toth}
\end{equation}
with
\begin{equation}
  n_s(T) = n_{s0}\exp\!\left[\chi\!\left(1-\frac{T}{T_0}\right)\right],\quad
  b(T)   = b_0\exp\!\left[\frac{\Delta H_0}{RT_0}\!\left(\frac{T_0}{T}-1\right)\right],\quad
  t(T)   = t_0 + \alpha\!\left(1-\frac{T_0}{T}\right),
  \label{eq:toth-T}
\end{equation}
where $p = c R_g T$ is the CO$_2$ partial pressure and $T_0 = 298.15$~K.
When $t_0\to1$, \eqref{eq:toth} recovers the Langmuir isotherm; for
$t_0\in(0,1)$ the isotherm is \emph{favorable} (concave downward), producing
self-sharpening wave behaviour.

Reference parameters fitted to amine-grafted $\gamma$-alumina pellets
\cite{Stampi2024,Grossmann2023}, used as the benchmark in Section~\ref{sec:gates}:
\begin{equation}
  n_{s0}=1.23~\text{mol\,kg}^{-1},\quad
  b_0=4839~\text{kPa}^{-1},\quad
  t_0=0.25,\quad
  \Delta H_0=70~\text{kJ\,mol}^{-1},\quad
  \chi=0,\quad
  \alpha=0.11.
  \label{eq:params}
\end{equation}
At $p_0 = 400~\text{ppm}\times101.325~\text{kPa} = 0.04~\text{kPa}$
and $T_0$, equation~\eqref{eq:toth} gives
$b_0 p_0 = 193.6$, $(b_0 p_0)^{t_0} = 0.1936^{1/4}\times10^{t_0\log10}\approx3.73$,
and $q_0^* \approx 0.48~\text{mol\,kg}^{-1}$.
These parameters are calibrated to $\gamma$-alumina and serve as a bracket
for PEI@SiO$_2$ pending project-specific isotherm measurements; the
higher heat of adsorption reported for PEI@silica fibre
($\Delta H_0=210~\text{kJ\,mol}^{-1}$, Pang et al.\ 2024) lies outside the
physically expected chemisorption range of 60--90~kJ\,mol$^{-1}$ and should be
verified against primary calorimetric data before use.

\subsection{Non-isothermal energy balance}

For the non-isothermal screening case,
\begin{multline}
  \bigl[\varepsilon\rho_g c_{p,g}+(1-\varepsilon)\rho_p c_{p,s}\bigr]
  \frac{\partial T}{\partial t}
  + \varepsilon\rho_g c_{p,g}\,u\frac{\partial T}{\partial z}
  \\
  = \lambda_\text{eff}\frac{\partial^2 T}{\partial z^2}
    + (1-\varepsilon)\rho_p(-\Delta H_\text{ads})\frac{\partial\bar{q}}{\partial t}
    - \frac{4h_w}{d_c}(T-T_w).
  \label{eq:energy}
\end{multline}
The wall-loss term is retained by default; for the SUTD column
($d_c=8.2$~mm) the wall-to-bed-volume ratio is large enough that
neglecting it is not justified.  The adiabatic limit $h_w=0$ is
retained as an explicit switch, not the default.

\subsection{Initial and boundary conditions}

\begin{equation}
  c(z,0)=0,\quad q_1(z,0)=q_2(z,0)=0,\quad T(z,0)=T_0.
  \label{eq:IC}
\end{equation}
Danckwerts inlet and zero-gradient outlet conditions:
\begin{equation}
  \left.\!\left(uc - \varepsilon D_L\frac{\partial c}{\partial z}\right)\right|_{z=0^+}
    \!=uc_0,\qquad
  \left.\frac{\partial c}{\partial z}\right|_{z=L}=0;\quad
  \left.\frac{\partial T}{\partial z}\right|_{z=0}
    =\left.\frac{\partial T}{\partial z}\right|_{z=L}=0.
  \label{eq:BC}
\end{equation}

\subsection{Dimensionless groups}
\label{sec:dimless}

Table~\ref{tab:dimless} collects all groups with physical interpretations and
order-of-magnitude estimates for the SUTD column
($L=0.15$~m, $d_c=8.2$~mm, $u=0.10$~m\,s$^{-1}$,
$\varepsilon=0.40$, $\rho_p=800$~kg\,m$^{-3}$, $q_0^*\approx0.48$~mol\,kg$^{-1}$).

\begin{table}[ht]
\centering
\caption{Dimensionless groups.
  $\varepsilon_p$ (particle porosity $\approx0.71$) enters film/pore correlations
  but not $\varepsilon$ (bed voidage $\approx0.40$); the two must not be
  conflated.}
\label{tab:dimless}
\renewcommand{\arraystretch}{1.2}
\begin{tabular}{@{}lllcc@{}}
\toprule
Symbol & Definition & Physical meaning & Estimate & Required limit \\
\midrule
$\delta_1$ & $\varepsilon c_0/[(1-\varepsilon)\rho_p q_0^*]$
  & gas-to-solid capacity ratio & $\sim2\times10^{-5}$ & $\ll1$ \\
$\mathrm{Pe}$ & $uL/D_L$
  & P\'eclet number & $\sim150$ & $\gg1$ \\
$\mathrm{NTU}_1$ & $k_1(1-\varepsilon)L/(\varepsilon u)$
  & transfer units (fast sites) & $O(1)$--$10$ & $O(1)$ \\
$\kappa$ & $k_2/k_1$
  & kinetic ratio & $0.01$--$0.1$ & $\ll1$ \\
$\Lambda$ & $(1-\varepsilon)\rho_p q_0^*/(\varepsilon c_0) = 1/\delta_1$
  & retardation factor & $\sim5\times10^{4}$ & $\gg1$ \\
$\Phi$ & $(1-\varepsilon)\rho_p|\Delta H_\text{ads}|q_0^*/
         [\varepsilon\rho_g c_{p,g}T_0]$
  & thermal loading & $O(1)$--$10$ & full model \\
$\mathrm{Bi}_w$ & $h_w d_c/(4\lambda_\text{eff})$
  & wall Biot number & $O(1)$ & retain ($d_c<10$~mm) \\
\bottomrule
\end{tabular}
\end{table}

The three-level hierarchy $\delta_1\ll 1/\mathrm{Pe}\ll 1$ holds for all
physically realistic DAC operating conditions and is the foundation of the
reduction in Section~\ref{sec:wave}.

%============================================================================
\section{Travelling-wave reduction}
\label{sec:wave}
%============================================================================

\subsection{Non-dimensionalisation}

Define $\hat{c}=c/c_0$, $\hat{q}_i=q_i/q_0^*$, $\hat{z}=z/L$, and the slow
time $\hat{t}=tu/(\Lambda L)$ scaled by the total stoichiometric breakthrough
time $t_b=\Lambda L/u$.  Under Assumption~\ref{ass:iso}, equation~\eqref{eq:gas}
becomes
\begin{equation}
  \delta_1\frac{\partial\hat{c}}{\partial\hat{t}}
  + \frac{\partial\hat{c}}{\partial\hat{z}}
  = \frac{1}{\mathrm{Pe}}\frac{\partial^2\hat{c}}{\partial\hat{z}^2}
    - \frac{\partial\hat{Q}}{\partial\hat{t}},
  \label{eq:ndgas}
\end{equation}
where $\hat{Q}=\hat{q}_1+\hat{q}_2$.  At $\delta_1=0$ and
$1/\mathrm{Pe}=0$, \eqref{eq:ndgas} is a scalar conservation law
$\partial\hat{c}/\partial\hat{z}+\partial\hat{Q}/\partial\hat{t}=0$
supporting shocks; dispersion ($1/\mathrm{Pe}$) regularises the shock and
gas accumulation ($\delta_1$) introduces a small compressibility correction.

\subsection{Wave ansatz and Rankine--Hugoniot speed}

We seek a constant-pattern solution
$\hat{c}=\hat{c}(\xi)$, $\hat{q}_i=\hat{q}_i(\xi)$, with
$\xi=\hat{z}-\hat{v}\hat{t}$.
At leading order ($\delta_1=0$, $1/\mathrm{Pe}=0$), substituting into
\eqref{eq:ndgas} and integrating from $+\infty$ (where $\hat{c}=\hat{Q}=0$)
to arbitrary $\xi$ gives
\begin{equation}
  \hat{c} + \hat{v}\hat{Q} = 0.
  \label{eq:rh-integral}
\end{equation}
The wave-frame LDF equations follow from
$\partial/\partial\hat{t}=-\hat{v}\,d/d\xi$:
\begin{align}
  \hat{v}\frac{d\hat{q}_1}{d\xi}
    &= \mathrm{NTU}_1\bigl(\eta\,\hat{q}^*(\hat{c}) - \hat{q}_1\bigr),
  \label{eq:waveq1}\\
  \hat{v}\frac{d\hat{q}_2}{d\xi}
    &= \kappa\,\mathrm{NTU}_1\bigl((1-\eta)\,\hat{q}^*(\hat{c}) - \hat{q}_2\bigr),
  \label{eq:waveq2}
\end{align}
where $\hat{q}^*(\hat{c})=q^*(c_0\hat{c},T_0)/q_0^*$ is the normalised
Toth isotherm.

\begin{theorem}[Dual-kinetic travelling wave]\label{thm:wave}
  Under Assumptions~\ref{ass:plug}--\ref{ass:iso}, with the Toth isotherm
  \eqref{eq:toth} favorable ($t_0\in(0,1)$) and the parameter ordering
  $0\leq\kappa<1$, $\delta_1\ll1$, $1/\mathrm{Pe}\ll1$:
  \begin{enumerate}
    \item[\emph{(i)}] \emph{(Wave existence.)}  There exists a unique
      travelling-wave solution with wave speed
      \begin{equation}
        v_{f,1} = \frac{uc_0}{(1-\varepsilon)\rho_p\,\eta\, q_0^*+\varepsilon c_0}
                 \approx \frac{uc_0}{(1-\varepsilon)\rho_p\,\eta\,q_0^*}
                 = \frac{v_f}{\eta},
        \label{eq:vf1}
      \end{equation}
      where $v_f = uc_0/[(1-\varepsilon)\rho_p q_0^*+\varepsilon c_0]$ is the
      Rankine--Hugoniot speed for the full equilibrated system.
    \item[\emph{(ii)}] \emph{(Perturbative slow-site correction.)}  For
      $0<\kappa\ll1$, the slow-site loading $\hat{q}_2(\xi)$ is bounded,
      $0\leq\hat{q}_2\leq 1-\eta$, and the induced correction to $\hat{c}$
      is $O(\kappa)$ uniformly in $\xi$.  In the time domain at $z=L$, the
      correction decays as $e^{-\kappa k_1(t-t_{b,1})}$ where
      $t_{b,1}=\eta\Lambda L/u$.
  \end{enumerate}
\end{theorem}

\begin{proof}
\textit{Step~1 ($\kappa=0$).}  Setting $\kappa=0$ in \eqref{eq:waveq2} gives
$\hat{q}_2\equiv0$, so $\hat{Q}=\hat{q}_1$.  The conservation integral
\eqref{eq:rh-integral} becomes $\hat{c}=-\hat{v}\hat{q}_1$.
Applying the upstream boundary condition
$\hat{c}\to1$, $\hat{q}_1\to\eta$ as $\xi\to-\infty$ gives
$1=-\hat{v}\eta$, hence $\hat{v}=-1/\eta$, recovering \eqref{eq:vf1}.

Substituting $\hat{c}=\hat{q}_1/\eta$ into \eqref{eq:waveq1}:
\[
  -\frac{1}{\eta}\frac{d\hat{q}_1}{d\xi}
    = \mathrm{NTU}_1\!\left(\eta\,\hat{q}^*\!\left(\frac{\hat{q}_1}{\eta}\right)
                              - \hat{q}_1\right).
\]
This is an autonomous scalar ODE on $[0,\eta]$.  At $\hat{q}_1=0$:
$\hat{q}^*(0)=0$, so the right-hand side vanishes.  At $\hat{q}_1=\eta$:
$\hat{q}^*(1)=1$ and the right-hand side equals $\mathrm{NTU}_1(\eta-\eta)=0$.
For the favorable Toth isotherm ($t_0\in(0,1)$), the normalised isotherm
satisfies $\hat{q}^*(\hat{c})>\hat{c}$ for all $\hat{c}\in(0,1)$
(the isotherm lies above the diagonal), so for $\hat{q}_1\in(0,\eta)$ the
right-hand side is strictly positive.  The ODE therefore has a unique
heteroclinic orbit connecting $\hat{q}_1=0$ (at $\xi\to+\infty$) to
$\hat{q}_1=\eta$ (at $\xi\to-\infty$), establishing the existence and
uniqueness of the wave \cite{Rhee1970}.

\textit{Step~2 (restoring $0<\kappa\ll1$).}
Write $\hat{Q}=\hat{q}_1+\hat{q}_2$.  The integral \eqref{eq:rh-integral}
gives $\hat{c}=-\hat{v}(\hat{q}_1+\hat{q}_2)$.  Since $0\leq\hat{q}_2\leq1-\eta$,
the correction $\hat{v}\hat{q}_2$ to $\hat{c}$ is bounded by
$\kappa\,\mathrm{NTU}_1\times O(1)$.  By standard continuous dependence on
parameters, the wave persists for small $\kappa>0$.

For $\hat{q}_2$: in the region behind the fast front ($\xi\ll0$, $\hat{c}\approx1$,
$\hat{q}^*\approx1$), equation~\eqref{eq:waveq2} reads
$(-1/\eta)\,d\hat{q}_2/d\xi=\kappa\,\mathrm{NTU}_1((1-\eta)-\hat{q}_2)$,
with solution $\hat{q}_2(\xi)=(1-\eta)(1-e^{\kappa\,\eta\,\mathrm{NTU}_1\xi})$.
Since $\xi<0$, this expression is bounded in $(0,1-\eta)$ and approaches
$1-\eta$ as $\xi\to-\infty$.  In the time domain at $z=L$, the
argument $\xi = (1 - t/t_{b,1})/(\Lambda/\eta)$ gives
$\hat{q}_2(t)\propto e^{-\kappa k_1(t-t_{b,1})}$, which decays to zero as
$t\to\infty$, confirming part~(ii).~$\square$
\end{proof}

\begin{remark}
  Theorem~\ref{thm:wave} requires a favorable isotherm (concave $q^*$,
  genuine nonlinearity) for the heteroclinic orbit to exist.  A linear Henry
  isotherm ($t_0=1$, $b_0\to0$) gives no heteroclinic orbit; the profile
  spreads as $t^{-1/2}$ (Klinkenberg regime), which is the Gate~A linear
  benchmark defined in Section~\ref{sec:gates}.
\end{remark}

\subsection{Hybrid outlet concentration}

Evaluating the wave at the column exit $z=L$ yields the breakthrough curve.
Let $t_{b,1}=\eta\Lambda L/u$ be the fast-site stoichiometric time.
In the DAC limit $y_0=c_0R_gT_0/p_a\to0$ (incompressible approximation
$\delta_1\to0$, valid to relative order $y_0\approx4\times10^{-4}$):
\begin{equation}
  \boxed{
  \frac{c(L,t)}{c_0}
  \;\approx\;
  \underbrace{\Bigl(1-e^{-k_1(t-t_{b,1})}\Bigr)}_{\text{fast-site rise}}
  \;-\;
  \underbrace{\frac{1-\eta}{\eta}\;\kappa\;e^{-\kappa k_1(t-t_{b,1})}}_{\text{slow-site tail}},
  \qquad t>t_{b,1},}
  \label{eq:hybrid}
\end{equation}
where $t_{b,1}=\eta(1-\varepsilon)\rho_p q_0^* L/(uc_0)$.
The first term is the Myers--Font (2020) exponential for the fast-site front
\cite{MyersFont2020}; the second is the $O(\kappa)$ tail correction derived in
Theorem~\ref{thm:wave}.  For finite $y_0$, the first term is replaced by the
exact Myers--Font compressible form
$(1-e^{-k_1(t-t_{b,1})})/(1-y_0 e^{-k_1(t-t_{b,1})})$.
The large-NTU Klinkenberg error-function front and the logistic approximation
are both recovered as special cases of the wave profile $\hat{c}(\xi)$ under
linear and Langmuir isotherms, respectively; the Toth wave requires numerical
integration of the wave-frame ODE.

\begin{corollary}[Homogeneous limit]\label{cor:hom}
When $\eta\to1$, $t_{b,1}\to t_b=(1-\varepsilon)\rho_p q_0^* L/(uc_0)$
and the tail term vanishes; \eqref{eq:hybrid} reduces to
$c(L,t)/c_0 = 1-e^{-k_1(t-t_b)}$, the Myers--Font single-LDF result.
\end{corollary}

\begin{corollary}[Sharp-front limit]\label{cor:sharp}
When $k_1\to\infty$ with $k_2/k_1=\kappa$ fixed, the fast-site rise
becomes a step at $t_{b,1}$ and the outlet curve is
$c/c_0 = 1 - ((1-\eta)/\eta)\kappa e^{-\kappa k_1(t-t_{b,1})}$
for $t>t_{b,1}$; the tail area $\int_{t_{b,1}}^\infty(1-c/c_0)\,dt$
equals $(1-\eta)t_{b,1}/\eta$.
\end{corollary}

\subsection{Asymmetry index and kinetic-model selection}

Define the asymmetry index of any measured breakthrough curve:
\begin{equation}
  \mathcal{A} = \frac{\tau_{95}-\tau_{50}}{\tau_{50}-\tau_5},
  \label{eq:asym}
\end{equation}
where $\tau_x$ is the time at which $c(L,\tau_x)/c_0=x/100$.
The Bohart--Adams, Thomas, and Yoon--Nelson models are all equivalent to
the logistic $c/c_0=[1+e^{k(\tau-t)}]^{-1}$ \cite{Hu2024}, which gives
$\mathcal{A}=1$ identically.
From Corollary~\ref{cor:hom}, $\mathcal{A}=1$ when $\eta=1$.
Cabrera-Codony et al.\ (2026) report $\mathcal{A}\approx5.4$ for PEI-fumed
silica at DAC concentrations, consistent with $\kappa\ll1$ and
$(1-\eta)/\eta\sim O(1)$ in \eqref{eq:hybrid}.
The practical selection rule is:
\begin{equation}
  \mathcal{A}\lesssim2 \;\Rightarrow\; \text{single-LDF (PFO) sufficient};\qquad
  \mathcal{A}\gtrsim3 \;\Rightarrow\; \text{escalate to dual-kinetic.}
  \label{eq:rule}
\end{equation}
Fitting proceeds in two stages: $k_1$ and $D_L$ are fit to the initial rise
(first $\sim70\%$ of the curve), then $\kappa$ and $\eta$ to the tail \cite{Stampi2024}.

%============================================================================
\section{Validation gates and parametric framework}
\label{sec:gates}
%============================================================================

\subsection{Three escalating validation gates}

Three gates must clear \emph{in order}; skipping a gate is not permitted.
Each tests a distinct physical ingredient.

\medskip
\textbf{Gate~A (Week~4) --- Klinkenberg linear benchmark.}
Set $q^*=K_H c$ (Henry isotherm, $\kappa=0$).
The analytical Klinkenberg solution \cite{Hu2024} is
\begin{equation}
  \frac{c(L,t)}{c_0}
  = \tfrac{1}{2}\Bigl[1+\mathrm{erf}\!\bigl(\sqrt{\zeta}-\sqrt{\tau}\bigr)\Bigr]
  + O(\zeta^{-1}),
  \label{eq:klink}
\end{equation}
valid for $\zeta\geq2$, where
$\zeta=k_H a_p(1-\varepsilon)L/(\varepsilon u)$ and
$\tau=k_H a_p K_H(1-\varepsilon)(t-\varepsilon L/u)/(\varepsilon u)$
\cite{Hu2024}.
The Method-of-Lines solver must reproduce \eqref{eq:klink} with cumulative
mass-balance error below 1\%.

\medskip
\textbf{Gate~B (Weeks~5--6) --- Rankine--Hugoniot front velocity.}
With the Toth isotherm \eqref{eq:toth}--\eqref{eq:params} active,
the simulated front velocity $v_\text{sim}=L/\tau_{50}$ must satisfy
$|v_\text{sim}/v_{f,1}-1|\leq0.15$.
This gate verifies both the nonlinear isotherm implementation and the
spatial convergence of the solver.  The tolerance is fixed at $\pm15\%$;
a stretch target of $\pm10\%$ may be reported separately.

\medskip
\textbf{Gate~C (Week~6) --- Stampi-Bombelli et al.\ (2024) breakthrough time.}
With parameters~\eqref{eq:params} scaled to the column geometry of
\cite{Stampi2024}, the simulated $\tau_{BT}$ (defined at $c/c_0=0.05$) must
fall within $\pm20\%$ of the reported experimental value.  Only after
Gate~C clears may parametric sweeps begin.

\subsection{Parametric hypotheses}

The five falsifiable hypotheses governing the OAT sweep are stated with
pre-committed acceptance thresholds:

\begin{description}
  \item[H1.] Increasing $u$ shortens $\tau_{BT}$ and widens the MTZ;
    the $\tau_{BT}$--$u$ slope matches $v_{f,1}$  within $\pm15\%$.
  \item[H2.] Increasing $c_0$ raises dynamic capacity $q_\text{dyn}$ and
    reduces $\tau_{BT}$; Toth predicts $q_\text{dyn}$ within $\pm20\%$ of
    the measured value.
  \item[H3.] $\tau_{BT}$ scales linearly with $L$ within $\pm10\%$ in the
    constant-pattern regime (NTU$_1>3$).
  \item[H4.] Raising $T_\text{ads}$ reduces $q_\text{dyn}$ through
    the van~'t~Hoff dependence of $b(T)$ in \eqref{eq:toth-T}.
  \item[H5.] The simulated shock velocity matches \eqref{eq:vf1} within
    $\pm15\%$ whenever $\mathrm{NTU}_1>5$.
\end{description}

The nine-run OAT design has one baseline ($u=0.10$~m\,s$^{-1}$,
$c_0=400$~ppm, $L=0.15$~m, $T=25$~°C) and two levels per parameter.
Results are visualised on a (Pe, NTU$_1$) regime map: Pe varies with $u$
and $L$; NTU$_1$ varies with $k_1$, $L$, and $u$.  The map separates the
sharp-shock regime (NTU$_1\gg1$, H5 territory), the kinetic-limited regime
(NTU$_1\sim1$, H1--H4 territory), and the dispersion-dominated regime
(Pe$\ll1$).

%============================================================================
\section{Discussion}
\label{sec:resolution}
%============================================================================

\subsection{Notation conflicts resolved}

Several conflicts between sources are standardised here to prevent
inconsistent numerical implementation.

\textbf{LDF driving force.}
The form $\partial\bar{q}/\partial t = k(q^*-\bar{q})$
(solid-loading driving force, rate constant in s$^{-1}$) is used throughout.
The alternative form $\partial\bar{q}/\partial t = k'(c-c^*)$
(gas-concentration driving force) uses $k'=k\,dq^*/dc$; the two coefficients
are numerically different by a factor equal to the local isotherm slope and
must not be exchanged.

\textbf{Gas-to-solid capacity ratio.}
The dimensionless group $\delta_1 = \varepsilon c_0/[(1-\varepsilon)\rho_p q_0^*]$
defined in Table~\ref{tab:dimless} is the gas-to-solid capacity ratio.
A distinct group $Lk_q/u_0\approx0.027$ (reaction-to-advection time ratio)
appearing in some sources carries different physical meaning and must not
be labelled $\delta_1$.

\textbf{Void fractions.}
$\varepsilon$ (bed voidage, $\approx0.40$) and $\varepsilon_p$ (particle porosity,
$\approx0.71$ for $\gamma$-alumina \cite{Stampi2024}) appear in separate
sub-models (column balance vs.\ intraparticle diffusion) and must not
be conflated.

\subsection{Scope and limitations}

The analysis covers single-component CO$_2$/N$_2$ adsorption, with
non-isothermal effects entering at order $\Phi\cdot\delta_1$ (small but
included in the numerical model), constant superficial velocity (Ergun
pressure drop negligible at the velocities considered), and no humidity.
Humidity enhancement of CO$_2$ uptake on PEI sorbents is reported in some
studies but remains mechanistically contested; it is excluded from the
present scope.

Whether PEI@SiO$_2$ granules require the dual-kinetic closure will be
determined from the asymmetry index $\mathcal{A}$ measured on the first
experimental curve.  Stampi-Bombelli et al.\ \cite{Stampi2024} found that
the dual-kinetic model collapses to single-LDF for packed beds of pellets
(contrasting with their monolith, where $\kappa\approx0.01$ is essential).
If $\mathcal{A}\lesssim2$, the single-LDF model (Corollary~\ref{cor:hom})
is sufficient and the theoretical results of this paper reduce to those of
Myers \& Font \cite{MyersFont2020}.

\subsection{Empirical models as failed-comparison cases}

Hu et al.\ \cite{Hu2024} prove that Bohart--Adams, Thomas, and Yoon--Nelson
are algebraically equivalent to the logistic $c/c_0=[1+e^{k(\tau-t)}]^{-1}$,
so fitting all three to the same dataset tests nothing.  Their rate
constants also drift with operating conditions (a threefold increase in
inlet concentration nearly doubles $k_{YN}$ for toluene on activated
carbon \cite{Hu2024}), confirming that these are shape-fitting parameters,
not physical constants.  Accordingly, Bohart--Adams, Thomas, and Yoon--Nelson
appear in this study only as failed-comparison cases; the benchmarks are the
analytically grounded Gates~A--C.

%--- References -------------------------------------------------------------
\begin{thebibliography}{99}

\bibitem{IEA2022}
IEA (2022). \textit{Direct Air Capture 2022}.  International Energy Agency,
Paris.

\bibitem{IPCC2023}
IPCC (2023). \textit{Climate Change 2023: Synthesis Report}.  IPCC, Geneva.

\bibitem{Pedrozo2025}
Pedrozo, H.\,A., Br\"uggemann, T., Serna, L.\,Y., Tsitsonis, V., Ajenifuja, A.,
  Sch\"afer, P., Kearns, J., \& Mazzotti, M.\ (2025).
Optimization of direct air capture processes using reactive transport models
  of adsorption--desorption cycles.
\textit{Computers \& Chemical Engineering}, 204, 109379.

\bibitem{Bollini2012}
Bollini, P., Choi, S., Drese, J.\,H., \& Jones, C.\,W.\ (2012).
Oxidative degradation of aminosilica adsorbents relevant to postcombustion
  CO$_2$ capture.
\textit{Energy \& Fuels}, 25(5), 2416--2425.

\bibitem{Kalyanaraman2015}
Kalyanaraman, J., Fan, Y., Labreche, Y., Lively, R.\,P., Kawajiri, Y.,
  \& Realff, M.\,J.\ (2015).
Modeling and experimental validation of carbon dioxide sorption on hollow
  fibers loaded with silica-supported poly(ethylenimine).
\textit{Chemical Engineering Journal}, 259, 737--751.

\bibitem{Ohs2018}
Ohs, B., Kr\"{o}del, M., \& Wessling, M.\ (2018).
Adsorption of carbon dioxide on solid amine-functionalized sorbents:
  A dual kinetic model.
\textit{Separation and Purification Technology}, 204, 13--20.

\bibitem{Hu2024}
Hu, Q., Li, J., Liu, Z., Sheng, T., Shen, J., \& Xie, J.\ (2024).
A critical review of breakthrough models with analytical solutions in a
  fixed-bed column.
\textit{Journal of Water Process Engineering}, 59, 105065.

\bibitem{Stampi2024}
Stampi-Bombelli, V., Storione, A., Grossmann, Q., \& Mazzotti, M.\ (2024).
On comparing packed beds and monoliths for CO$_2$ capture from air through
  experiments, theory, and modeling.
\textit{Industrial \& Engineering Chemistry Research}, 63(26), 11637--11653.

\bibitem{Grossmann2023}
Grossmann, Q., Stampi-Bombelli, V., \& Mazzotti, M.\ (2023).
Developing versatile contactors for direct air capture of CO$_2$ through
  amine grafting onto alumina pellets and alumina wash-coated monoliths.
\textit{Industrial \& Engineering Chemistry Research}, 62(33), 13594--13611.

\bibitem{MyersFont2020}
Myers, T.\,G., \& Font, F.\ (2020).
Mass transfer from a fluid flowing through a porous media.
\textit{International Journal of Heat and Mass Transfer}, 163, 120374.
[\textit{Note:} the article number 120374 appears also for an unrelated
reference in the literature; the primary PDF should be verified before
final submission.]

\bibitem{Rhee1970}
Rhee, H.-K., Aris, R., \& Amundson, N.\,R.\ (1970).
On the theory of multicomponent chromatography.
\textit{Philosophical Transactions of the Royal Society A}, 267, 419--455.

\bibitem{Glueckauf1955}
Glueckauf, E.\ (1955).
Theory of chromatography, Part~10: Formula\ae{} for diffusion into spheres
  and their application to chromatography.
\textit{Transactions of the Faraday Society}, 51, 1540--1551.

\end{thebibliography}

\end{document}